% %%%%%%%%%%%%%%%%%%%%%%%%%%%%%%%%%%%%%%%%%%%%%%%%%%%%%%%%%%%%%%%%%%%%%
% %% State-of-the-Art %%
% InnerSP
% %%%%%%%%%%%%%%%%%%%%%%%%%%%%%%%%%%%%%%%%%%%%%%%%%%%%%%%%%%%%%%%%%%%%%

\subsection{InnerSP}
\label{sec:soa:innersp}

InnerSP~\cite{Baek2021_InnerSP} is a SpMSpM accelerator based on the \algo{Gust} algorithm\footnote{The authors call this \emph{row-wise inner-product}, or confusingly sometimes \emph{inner-product}.}. The authors of InnerSP justify their choice by showing quantitatively that with the \algo{OP} the \emph{POs} represent a large volume of data, and that for many matrices the \emph{reductions} can not be done fully on-chip and thus must be buffered in main memory. Normally, with \algo{Gust} algorithm, $B$ must be read repeatedly, but InnerSP proposes a special cache for $B$ which exploits the spatial locality.

The architecture of InnerSP is shown in Figure~\ref{fig:soa:innersp}. The \emph{A reader} reads the $A$-rows. Two separate caches are used for $B$, one that stores the row pointers and the other which stores column indices and values. The \emph{B reader} reads the required $B$-rows, hopefully getting the data from the caches. Parallel multipliers perform the multiplications and then a hash-unit generates a hash based on the indices in $C$. A \emph{reduction} unit based on hash tables contains parallel comparators which search for matching hash values and performs the \emph {reductions}. When all the operations for a $A$-row are completed, the \emph{C writer} writes the $C$-row back to main memory.

% Figure located in Accelerators.tex

The architecture is efficient, as long as the \emph{Hash Reduction Unit} does not overflow. When overflows occur, the extra entries are temporarily stored in main memory, but there is a significant performance hit. To minimize the chance of overflows, InnerSP uses a pre-scanner to estimate the number of non-zero entries in the upcoming $C$-rows. The estimate primarily considers the $NNZ$ in the $A$-rows. If the estimated number of non-zero entries in $C$ exceeds the capacity of the \emph{Hash Reduction Unit}, then the row is split and processed in two passes. Conversely, if the number of non-zero entries in $C$ is small, two $A$-rows can be processed simultaneously.

To improve the performance of the $B$-cache, the authors exploit a technique from~\cite{Balaji2021_Popt}. The non-zero elements in the $A$-rows provide direct information about which $B$-rows are required. The column index table of $A$ is already pre-fetched by the \emph{A reader}, and when a row must be evicted from the $B$-cache, it chooses the ones which will not be required based on the upcoming entries in the column index table of $A$. This technique (also used by SpArch~\cite{Zhekai2020_SpArch}) maximizes the reuse of $B$-rows.

InnerSP is an accelerator based on \algo{Gust} algorithm which achieves high reuse of the $B$-rows, through a look-ahead in the $A$-\emph{metadata}.